ODE 的未知量是一条曲线，一个时刻对应一个数。温度场、弦的形状、流体的速度不是这样：每个时刻要交出的是一整条剖面、一整张曲面，等式里时间导数与空间导数同时出场。这就是偏微分方程。Part 15 前面建立的概念——初值问题、稳定性、能量、谱——在这里几乎都有对应物，但每一样都要重新检查一遍，因为多出来的空间维度带来了 ODE 里没有的东西：边界。

边界不是附加说明。同一个热方程，一端贴冰袋与一端裹保温层，解完全不同；连解存不存在、唯一不唯一，都要看边界条件给得对不对。方程管局部怎样演化，初值管从哪里出发，边界管在边缘处怎样与外界交换，三样凑齐才谈得上唯一解。

\begin{center}
\begin{tikzpicture}[>=stealth]
  \node[draw] (a) at (0,0) {\small 未知量是场};
  \node[draw] (b) at (3.7,0) {\small 抹平 \(\cdot\) 搬运 \(\cdot\) 平衡};
  \node[draw] (c) at (8.2,0) {\small 存在 \(\cdot\) 唯一 \(\cdot\) 连续依赖};
  \node[draw] (d) at (2.4,-1.9) {\small 分离变量：一族 ODE};
  \node[draw] (e) at (7.4,-1.9) {\small 弱解与半群};
  \draw[->] (a) -- (b);
  \draw[->] (b) -- (c);
  \draw[->] (c.south) -- ++(0,-0.85) -| (d.north);
  \draw[->] (d) -- (e);
\end{tikzpicture}
\end{center}

\emph{本章路线：先认清未知量成了场，再按三类行为分家，查一遍适定性，然后动手求解，最后放宽``解''的定义。}

\subsection{三类行为是这一章最该带走的判断}

热方程抹平：高频按 \(e^{-\kappa k^2t}\) 衰减，越细的结构死得越快，再粗糙的初值一瞬间也变成无限可微，代价是被抹掉的信息拿不回来。波动方程搬运：形状沿特征线以有限速度平移，尖角还是尖角，正则性既不增也不减。Laplace 方程平衡：没有时间，内部不允许出现局部极值，边界上的值一次性锁死整个内部。

认出手上的方程属于哪一类，就知道该配什么数据、解会有多光滑、扰动往哪个方向传，也知道数值上该防什么——抛物型的显式格式受 \(\Delta t\lesssim\Delta x^2\) 的刚性限制，双曲型受 CFL 条件约束，椭圆型没有时间步可言，代价转移到一次大型稀疏求解的条件数上。

\subsection{适定性的第三条最容易被跳过}

存在、唯一、对数据连续依赖。前两条读起来天经地义，第三条决定问题能不能交付。倒着解热方程时高频以 \(e^{+\kappa k^2\abs{t}}\) 放大，观测里那点传感器噪声会被推到台前，答案面目全非。去模糊、CT 重建、地震反演遇到的困难都是这一件事，用数值分析的话说，就是条件数没有上界。它属于问题本身，换求解器救不回来，只能靠正则化改问题。

\subsection{怎么读这一章}

首次阅读以前三篇为主：从``未知量是场''理解维数为何骤增，从热、波、Laplace 三个原型预判信息怎样传播，再用三条标准检查手上的问题是否适定。第四篇是规则区域上的计算样板，顺带解释了 Fourier 展开的来历——它是分离变量在最对称几何上的产物；第五篇是严格化的入口，需要弱解、有限元或无限维演化时回读。有一处提醒：椭圆、抛物、双曲的分类首先是二阶线性方程主部的局部分类，不要机械地套到任意非线性方程上，变系数问题的类型可以在区域内部改变。

\subsection{本章篇目}

以下篇目按概念依赖排列；若从具体问题进入，也可以先打开最接近当前困惑的一篇，再沿篇内链接回补。

\begin{itemize}
\tightlist
\item \hyperref[sec:15-Differential-Equations-and-Dynamical-Systems/03-PDE-Introduction/01]{``从 ODE 到 PDE：未知量从一条曲线变成一场''}；
\item \hyperref[sec:15-Differential-Equations-and-Dynamical-Systems/03-PDE-Introduction/02]{``三类典型方程行为迥异''}；
\item \hyperref[sec:15-Differential-Equations-and-Dynamical-Systems/03-PDE-Introduction/03]{``适定性与能量方法''}；
\item \hyperref[sec:15-Differential-Equations-and-Dynamical-Systems/03-PDE-Introduction/04]{``分离变量：把 PDE 还原为一族 ODE''}；
\item \hyperref[sec:15-Differential-Equations-and-Dynamical-Systems/03-PDE-Introduction/05]{``Sobolev 空间与半群：弱解的严格框架''}。
\end{itemize}

读完这五篇，控制部分的 HJB 方程、数值分析部分的谱方法与边值问题、扩散模型里的前向加噪，都可以借这里的分类与能量语言来读。
